% 第二章：相关工作
\chapter{相关工作}

\section{视觉模型}
视觉模型的发展经历了从“以视觉编码器为核心的表征学习”到“视觉-语言对齐的基础模型”，再到“面向复杂任务的多模态指令跟随与推理”的演进。本文的研究问题（多模态数据合成框架）与后续两条任务线（GUI、世界知识）都建立在这一演进脉络之上，因此本节重点回顾与本文最相关的视觉-语言模型（VLM）与多模态大模型（MLLM）的关键路线。

\subsection{视觉表征与视觉-语言对齐}
视觉表征学习是多模态模型的基础。早期的视觉模型主要依赖卷积神经网络（CNN）进行特征提取，其中以残差网络（ResNet）为代表的架构通过更深网络与稳定训练显著提升了表征能力\citep{he2016resnet}。随后，Vision Transformer (ViT) 将Transformer架构引入视觉领域，通过将图像切分为Patch序列进行处理，强化了对全局上下文的建模能力，并推动了“以预训练规模驱动性能”的范式\citep{dosovitskiy2021vit}。

在此基础上，对比学习（Contrastive Learning）成为视觉与语言建立共享语义空间的重要途径。CLIP通过海量图文对进行对比学习，使模型获得了强迁移的视觉表示能力，并在零样本分类等任务上表现突出\citep{radford2021learning}。与CLIP相近的路线还包括ALIGN等工作，强调“弱噪声文本监督+更大规模数据”的可扩展性，从而进一步巩固了开放词表视觉语义对齐的可行性\citep{jia2021align}。这类方法的关键贡献在于将视觉表征从“固定类别监督”转向“语言可描述的开放世界概念”，为后续多模态生成式建模提供了通用的语义对齐基础；同时也提示了一个事实：模型能力上限往往由数据覆盖度与噪声控制能力共同决定。

\subsection{从VLM到MLLM：融合大语言模型的多模态生成}
随着大语言模型（LLM）能力的爆发，主流路线逐渐从“检索/匹配式VLM”走向“以LLM为核心的生成式MLLM”。这一范式利用视觉编码器提取视觉特征，通过适配器（Adapter）映射到LLM的词嵌入空间，最终由LLM统一完成理解、生成与推理。

Flamingo采用了在冻结的LLM中插入交叉注意力层（Cross-Attention）的架构，并在海量交错图文数据上进行训练，展示了在少样本条件下的强大通用能力\citep{alayrac2022flamingo}。BLIP-2进一步提出了Q-Former结构，通过一组可学习的Query向量与图像编码器交互，实现了更高效的跨模态桥接，在保持训练效率的同时获得了强泛化能力\citep{li2023blip2}。

在“指令跟随（Instruction Following）”框架下，多模态模型通过视觉指令微调（Visual Instruction Tuning）进一步获得了更好的对话式交互能力。LLaVA系列工作系统化探索了视觉指令微调的数据构造与训练策略，证明了简单的线性映射或MLP即可实现高效的模态对齐\citep{liu2024visual,liu2023improved}。与此并行的路线还包括：基于“统一理解与生成”的BLIP框架\citep{li2022blip}，以及面向指令调优的InstructBLIP等工作\citep{dai2023instructblip}，它们都强调通过更结构化的监督信号让模型输出更可控、更贴近真实任务分布。

当训练规模进一步扩大时，PaLI等工作展示了“语言与图像共同扩展”的另一条路径，即通过多语种与多任务的联合扩展获得更强泛化\citep{chen2022pali}。此外，GPT-4o等原生多模态模型（Native Multimodal Models）的出现，进一步打破了模态间的界限，实现了端到端的跨模态理解与生成\citep{openai2024gpt4o}；Gemini等技术报告也表明，面向实际交互的多模态能力（如工具使用、长上下文、多步推理）正在成为系统级目标\citep{team2023gemini}。总体而言，这一阶段的核心变化是：模型能力的瓶颈逐渐从“是否能对齐视觉与语言”转向“是否具备足够高质量、可控结构的监督信号”，而数据构造与质控机制开始与模型架构同等重要。

近年来，先进MLLM进一步将能力边界推向“长上下文+视频+代理（agent）”的综合场景。一方面，Gemini系列与GPT-4o代表了闭源原生多模态系统在实时交互与端到端建模上的方向\citep{team2023gemini,openai2024gpt4o}；另一方面，开源侧也在快速演进，例如Qwen-3-VL系列通过更长上下文与更强的视频/交错多模态建模能力增强复杂任务覆盖\citep{qwen3vl2025}，而MiMo-VL等工作则探索将大规模预训练与强化学习式对齐结合，以在推理、GUI/agent等任务上获得更强的可执行性与鲁棒性\citep{mimovl2025}。

\subsection{视觉推理任务与结构化视觉理解}
除通用图文理解外，图表、文档、截图等“结构化视觉输入”对模型的视觉解析与逻辑推理提出了更高要求。这类数据通常包含密集的文本信息和复杂的布局结构，单纯的图像描述（Captioning）往往难以涵盖其全部语义。

ChartQA提出了包含视觉与逻辑推理的图表问答基准，推动模型从简单的“读图”走向复杂的“推断”\citep{masry2022chartqa}。DePlot创新性地将图表理解转化为“图到表”的结构化转换任务，从而以可执行的中间表征（表格）支撑后续的符号推理\citep{liu2022deplot}；MatCha进一步引入了数学推理与图表反渲染（De-rendering）机制来增强预训练效果\citep{liu2022matcha}。Pix2Struct则提出将“截图解析”作为预训练任务，强调对UI、网页等截图的结构化HTML/DOM树恢复能力\citep{lee2023pix2struct}。这些工作共同表明：当任务依赖结构化信息（布局、文本、关系）时，需要更强的结构化中间表示与针对性的数据构造策略。

\subsection{生成式视觉模型与可控图像合成}
在图像生成方向，扩散模型（Diffusion Models）推动了高分辨率、高保真生成能力的快速发展。以SDXL为代表的高质量扩散模型显著提升了高分辨率图像合成的质量与细节表现\citep{podell2023sdxl}，为“文本$\rightarrow$图像”的逆向数据工程提供了坚实的工具基础。对于本文的研究而言，生成模型的价值不仅在于生成图像本身，更在于为多模态数据合成提供可控的视觉证据（Visual Evidence），从而能够以低成本的方式弥补真实数据在长尾场景下的稀缺性以及图文弱对齐的问题。

\section{用户图形界面（GUI）}
面向GUI的视觉理解与交互是多模态模型落地的重要方向之一。与自然图像相比，GUI截图具有控件密集、文字细小、布局强结构化与状态易变等特征，使得模型需要同时具备细粒度视觉感知能力（小目标定位、OCR相关细节）、结构化理解能力（控件关系、层级与语义绑定）以及可执行输出能力（产生可解析的坐标或动作序列）。

\subsection{GUI视觉特征与理解挑战}
GUI图像与自然场景图像存在显著差异，其挑战不仅来自视觉观测本身，也来自“可执行交互”所要求的结构化与可控输出。首先，GUI界面通常包含大量细小的图标与高密度文本（按钮、菜单、表格、提示信息等），这要求模型具备更高分辨率的细粒度感知与稳定的OCR/文本检测能力；在工程实践中，轻量OCR系统（如PP-OCR）常被用作GUI理解链路中的关键组件\citep{du2020ppocr}。其次，GUI语义强依赖布局与层级关系（DOM树、View Hierarchy、可访问性树A11y等），因此仅基于图像级描述往往不足，需要“像素$\rightarrow$结构”的解析能力来恢复可执行的中间表示。Pix2Struct通过将截图解析作为预训练目标，验证了结构化表示对截图类输入的增益\citep{lee2023pix2struct}。

进一步地，GUI任务的核心输出通常不是自然语言描述，而是可执行动作（点击/输入/滚动等）及其对应的目标元素定位，这使得元素级grounding成为能力瓶颈之一。SeeClick等工作专门聚焦GUI点击与定位，强调更严格的坐标/区域监督与视觉指向能力\citep{liu2024seeclick}；OS-ATLAS则从跨平台角度提出统一动作空间与大规模定位语料，推动模型在不同操作系统与界面风格上的泛化\citep{wu2024osatlas}。

最后，真实GUI交互具有动态性与多步性：界面状态随操作而变化，且任务往往需要“观察$\rightarrow$规划$\rightarrow$执行$\rightarrow$纠错”的闭环。GUI-WORLD等数据集面向动态GUI内容与多场景任务构建评测，揭示了仅凭单帧理解的局限\citep{guiworld2024}；GUICourse提出从通用VLM到GUI智能体的分阶段数据与训练范式，强调OCR、GUI知识与交互能力的递进式培养\citep{guicourse2024}。同时，面向GUI智能体的系统性综述也指出：数据基础设施、评测协议与错误恢复机制正在成为影响落地可靠性的关键因素\citep{llmguiaentsurvey2024}。

\subsection{GUI数据集与基准构建}
高质量的数据集是推动GUI理解发展的关键。Rico数据集是早期最具代表性的移动端UI数据集，包含了大量Android应用截图及其对应的视图层级数据，为UI设计检索和布局生成提供了基础\citep{deka2017rico}。然而，原始的Rico数据存在大量噪声和缺失。Clay进一步提出了基于学习的方法来对原始UI布局进行去噪和修复，显著提升了数据的可用性\citep{li2022clay}。

近期，随着大模型对数据规模需求的增加，MobileViews构建了百万量级的移动端GUI数据集，涵盖了更广泛的应用场景和界面类型\citep{gao2024mobileviews}。此外，为了支持动态交互任务，DroidBot等自动化测试工具被广泛用于通过随机或策略性探索来收集动态的UI交互轨迹\citep{droidbot}。这些数据集的演进反映了从“静态布局分析”向“动态交互理解”的趋势。

\subsection{GUI智能体与动作模型}
在模型层面，研究者们致力于构建能够像人类一样操作设备的GUI智能体。Pix2Struct通过将截图解析任务作为预训练目标，证明了从像素中恢复结构化表示（如HTML）能够有效提升模型对UI的理解能力\citep{lee2023pix2struct}。

更进一步，OS-Atlas提出了通用的GUI动作模型（Action Model），通过在海量跨平台GUI数据上进行从粗粒度到细粒度的动作预测训练，实现了对不同操作系统（移动端、桌面端）的泛化控制能力\citep{wu2024osatlas}。这类工作的核心在于将视觉感知与动作空间（如点击、滑动、键盘输入）进行对齐，使模型不仅能“看懂”界面，还能“操作”界面。这一观察与本文第\ref{chap:gui}章采用的多源融合、结构化标注与质控迭代思路一致，即通过高质量的合成数据来弥补真实交互数据的稀缺性。

除了移动端与桌面端，Web场景为GUI智能体提供了更开放、可复现的评测环境。World of Bits（MiniWoB等）通过将网页交互任务抽象为一组可自动评测的微任务，推动了从“脚本/规则自动化”向“学习型交互策略”的转变\citep{shi2017worldofbits}。随着任务难度上升，Mind2Web等工作开始关注真实网页上的长程任务与多步决策，强调智能体在开放网页环境下的泛化与鲁棒性\citep{deng2023mind2web}；WebArena进一步将评测扩展到更贴近真实互联网的复杂环境，突出端到端评测中“检索、阅读、规划、执行、纠错”的闭环能力\citep{zhou2023webarena}。

与上述“任务级”评测相配套，界面元素的精确定位（例如“点击哪个按钮”）也逐渐成为能力瓶颈之一。SeeClick等工作将GUI控制进一步聚焦到细粒度点击与定位问题，提示模型需要更强的视觉指向（grounding）能力以及更严格的坐标/区域监督\citep{liu2024seeclick}。

\section{世界知识}
世界知识（World Knowledge）能力指模型对事实、概念与关系网络的掌握，以及在视觉证据与知识之间建立可验证链接并完成多步推理的能力。对于视觉语言模型而言，世界知识增强的难点通常不在于“语言侧是否存储过知识”，而在于“视觉侧是否能稳定锚定实体实例”，以及“推理链路是否可控、可迁移”。

\subsection{知识型视觉问答基准与信息检索需求}
世界知识能力在多模态研究中常通过“知识型视觉问答（Knowledge-based VQA）”体现：问题往往无法仅从图像像素直接回答，必须结合外部常识、百科或领域知识。OK-VQA提出了需要外部知识的VQA基准，推动研究从“看图问答”走向“看图+用知识回答”\citep{marino2019okvqa}；A-OKVQA进一步强化了对常识与知识组合推理的要求，使得模型需要在候选知识、视觉证据与推理链之间建立更稳定的对应关系\citep{schwenk2022aokvqa}。随着开放世界问题增多，WebQA等工作将多模态问答扩展到“跨模态多跳检索与推理”，强调证据可追溯性与多源证据融合\citep{chang2022webqa}；InfoSeek等基准则进一步强调“信息检索式”的问答设定，要求模型能够主动检索并在证据约束下回答，从而更适合评估事实性与幻觉控制能力\citep{chen2023infoseek}。

\subsection{知识推理与思维链（CoT）}
推理方法与对齐策略对世界知识问答质量至关重要。链式思维提示（Chain-of-Thought, CoT）表明显式推理过程可以显著提升LLM在复杂任务上的表现\citep{wei2022chain}。围绕CoT，后续研究进一步从\textbf{推理路径的可靠性}与\textbf{推理过程的可控性}两方面增强推理能力：Self-Consistency通过对同一问题采样多条推理链并进行一致性聚合，缓解单一路径的偶然错误\citep{wang2022selfconsistency}；Tree-of-Thoughts则将“思维”视为可搜索的中间状态，引入分支探索与回溯以提升需要规划/搜索任务的解题能力\citep{yao2023tot}。

更进一步，世界知识问题往往要求模型在外部证据约束下推理，单纯生成推理链容易产生幻觉或证据缺失。ReAct提出将“推理轨迹”与“行动（如检索/查询工具）”交错生成，通过与外部知识源交互来降低幻觉与错误传播\citep{yao2023react}。在多模态场景中，相关思想也被扩展为“多模态推理链”：Multimodal-CoT通过分阶段生成多模态依据与答案推断，提升知识型VQA等任务的稳定性\citep{zhang2023multimodalcot}；MM-REACT进一步将语言模型与多种视觉专家工具结合，实现“看图$\rightarrow$调用专家$\rightarrow$推理$\rightarrow$行动”的系统范式\citep{yang2023mmreact}；ViperGPT则通过生成并执行Python程序来组合视觉子模块，提高可解释性与可组合推理能力\citep{suris2023vipergpt}。这些工作共同表明：面向世界知识的多模态推理需要“实体锚定$\rightarrow$证据检索/调用$\rightarrow$可验证推理”的闭环，而CoT类方法为构造可迁移、可审计的推理监督提供了重要工具。

\subsection{指令微调与对齐策略}
为了让模型更好地遵循人类指令并输出高质量的回答，指令微调（Instruction Tuning）成为了关键技术。InstructGPT通过基于人类反馈的强化学习（RLHF）显著提升了模型遵循指令的能力\citep{ouyang2022training}。

进一步地，为了解决高质量指令数据稀缺的问题，Self-Instruct提出了通过模型自举生成指令数据以实现对齐的方法\citep{wang2023selfinstruct}。Orca等工作则进一步探索从强教师模型（如GPT-4）的“解释轨迹（Explanation Traces）”中进行渐进式学习，通过模仿教师模型的推理过程来提升学生模型的复杂推理能力\citep{mukherjee2023orca}。这些研究共同指出：当任务需要多步推断与事实一致性时，仅监督最终答案可能诱发“短路映射”或幻觉，过程性监督与可验证约束更利于能力迁移。

\subsection{事实性评估与多模态知识融合}
衡量与改进模型事实性也是世界知识增强的重要组成部分。SimpleQA聚焦短回答场景下的事实性评测，为分析模型在事实正确性上的薄弱点提供了参考\citep{wei2024simpleqa}。在方法层面，检索增强生成（Retrieval-Augmented Generation, RAG）等工作通过在生成前检索外部文档并在生成时融合证据，提供了一条提升事实性与可追溯性的通用路径\citep{lewis2020rag}。在多模态场景中，世界知识能力往往体现为“先识别再调用知识”的链路是否稳定、答案是否唯一可核验。本文第\ref{chap:world_knowledge}章即在这一思路下构造RecQA/KnowQA/FinalQA分解，并通过引入带CoT的FinalQA提升可迁移推理能力。

\section{多模态数据合成}
随着多模态大模型（MLLM）参数规模的不断增长，高质量数据的稀缺性逐渐成为限制模型性能提升的瓶颈。如何规模化构造高质量、强对齐且具备推理密度的数据成为当前研究的关键问题。现有工作主要围绕三个方向展开：提升Caption质量与对齐度、指令数据的进化与扩展、以及利用生成模型进行逆向数据工程。

\subsection{高质量Caption与视觉-语言对齐}
传统的图文数据集往往存在描述简短、信息遗漏等问题。ShareGPT4V通过利用GPT-4V生成详尽的图像描述（Dense Caption），证明了“更好的描述”可以显著改善视觉-语言对齐效果，进而提升下游任务的性能\citep{chen2023sharegpt4v}。这类工作的核心启示在于：Caption不仅是训练目标，更是连接视觉信号与语言语义的桥梁。通过将图像转化为可复用、可验证的文本证据，可以有效降低模态间的语义鸿沟。

从更宏观的角度看，数据规模与覆盖度同样决定了对齐的上限。LAION-5B等工作构建了超大规模开源图文对数据集，为训练与评测下一代图文对齐模型提供了基础设施\citep{schuhmann2022laion5b}；COYO-700M等数据集进一步展示了“更大规模弱监督数据”在训练可迁移视觉-语言表征上的价值\citep{byeon2022coyo}。与此同时，如何在规模扩展下控制噪声与偏差，也成为数据工程与模型训练共同面对的问题。

\subsection{指令进化与代理式数据生成}
在纯语言领域，Self-Instruct提出了通过模型自举（Bootstrapping）生成指令数据以实现对齐的方法，极大地降低了人工标注成本\citep{wang2023selfinstruct}。这一思想被迅速引入多模态领域。MMEvol提出了Evol-Instruct策略，通过多轮对话和指令改写来增强多模态指令的多样性与复杂度\citep{xu2024mmevol}。

更进一步，AgentInstruct引入了代理流（Agentic Flow）的概念，将数据生成过程组织为多步骤的代理协作流程，通过模拟用户与系统的交互来生成更具逻辑性和连贯性的数据\citep{mitra2024agentinstruct}。在结构化推理任务上，Synthesize Step-by-Step通过结合工具调用、模板填充与LLM生成，构建了包含推理步骤的图表问答数据，为“可执行/可验证”的数据构造提供了范例\citep{li2024synthesize}。这些方法共同体现了数据构造从“单轮生成”向“流程化、结构化生成”的演进趋势。

\subsection{合成数据幻觉与质量控制}
当数据构造从“人工标注”逐步转向“模型生成（Caption/指令/图像/轨迹）”时，合成数据不可避免地引入幻觉（Hallucination）与失真（Distortion）：模型可能在描述中虚构不存在的对象或属性、在指令中生成与界面不匹配的操作目标、或在多步轨迹中产生不可执行的状态转移。这些噪声一旦作为监督信号回灌训练，往往会以“错误对齐”的方式固化到模型中，并在规模扩张时被放大。

一类典型风险来自\textbf{Caption与对象一致性}：自动Caption器可能倾向于补全“常见但未出现”的对象，从而造成对象幻觉。POPE提出以投票式查询（Polling-based Query）更稳定地评测大视觉语言模型的对象幻觉问题，为合成文本的质量诊断提供了工具\citep{li2023pope}。在工程上，类似CLIPScore这类参考无关的图文一致性指标也常用于大规模过滤与重采样，以提升合成监督的可靠性\citep{hessel2021clipscore}。

另一类风险来自\textbf{指令-元素对齐与可执行性}：在GUI/交互数据合成中，若仅依赖模型“看图生成指令”，很容易产生目标元素不存在、边界框偏移或隐式指令理解失败等问题。UI-E2I-Synth通过引入大规模指令合成管道并聚焦元素比例失真、元素类型不平衡与隐式指令等难点，同时构建UI-I2E-Bench进行细粒度诊断，展示了“合成$\rightarrow$评测$\rightarrow$再合成”的质控闭环对降低幻觉的重要性\citep{liu2025uie2isynth}。

此外，合成数据还存在\textbf{递归训练的分布塌缩风险}：当模型反复使用自身或同类模型生成的数据进行训练时，长尾分布可能逐步消失并导致能力不可逆退化（Model Collapse）\citep{shumailov2023curse}。因此，当前多模态数据工程普遍强调：在扩大合成规模的同时，必须维持真实数据锚点、引入结构化验证器/可执行回放、并以多维指标持续监控噪声与偏差。本文后续提出的多源融合、结构化标注与质控迭代，也正是在这一“幻觉可诊断、质量可闭环”的原则下展开。

\subsection{从文本到图像的合成}
除了基于现有图像生成文本（Image-to-Text）的正向构造外，利用生成模型进行“文本到图像（Text-to-Image）”的逆向数据工程也日益受到关注。SynthVLM系统研究了合成图文数据集的构建策略，通过Caption清洗、扩散模型生成与自动质量筛选等环节，构建了大规模的高质量图文对\citep{liu2025synthvlm}。得益于SDXL等高质量扩散模型的发展，生成的图像在分辨率和语义一致性上已达到较高水平\citep{podell2023sdxl}。

对于本文所研究的“多模态数据合成框架”而言，正向与逆向两类数据流的结合，使得系统既能覆盖真实世界的长尾分布，又能在特定场景下进行可控的数据补齐，从而实现数据质量与规模的双重提升。

综合来看，相关工作为本文后续三部分内容提供了直接支撑：一类工作强调结构化锚点、精细监督与可执行对齐，对应第\ref{chap:gui}章中面向GUI理解的正向合成框架；一类工作强调caption中间表征、过滤、逆向数据组织与过程性推理监督，对应第\ref{chap:world_knowledge}章中面向世界知识的升级框架；而围绕指令进化、代理式生成、质量控制与正反向数据流的研究，则共同构成第5章Auto-Evol统一闭环框架的理论与方法基础。由此，本文并非将GUI、世界知识与统一框架视为三块彼此独立的内容，而是将其组织为一条逐步完善的数据合成框架演进路线。
